% 米氏散射：推導骨架（完整推導見 Mie (1908) 與 Bohren & Huffman (1983) 第四章）
本附錄只列出推導的骨架與本研究用到的結果，各步驟的完整推導見 \textcite{https://doi.org/10.1002/andp.19083300302} 的原始論文，符號與寫法依 \textcite{bohren1983Absorption} 第四章。

\subsubsection*{(一) 亥姆霍茲方程}
均勻、無源介質中的馬克士威方程式取時諧形式 $e^{-i\omega t}$ 後，電場與磁場皆滿足向量亥姆霍茲方程，且散度為零：
\begin{equation}
    \nabla^2\mathbf{E} + k^2 \mathbf{E} = 0, \qquad \nabla\cdot\mathbf{E}=0, \qquad k^2 = \omega^2 \epsilon \mu .
\end{equation}

\subsubsection*{(二) 球面邊界條件}
在球面 $r=a$ 上，電場與磁場的切向分量連續：
\begin{equation}
    \left(\mathbf{E}_{i} + \mathbf{E}_{s} - \mathbf{E}_{1}\right) \times \mathbf{e}_r = 0, \qquad \left(\mathbf{H}_{i} + \mathbf{H}_{s} - \mathbf{H}_{1}\right) \times \mathbf{e}_r = 0 ,
\end{equation}
其中下標 $i$、$s$、$1$ 分別代表入射場、散射場與球內場。

\subsubsection*{(三) 向量球面波函數}
由滿足純量亥姆霍茲方程的生成函數 $\psi$ 建構
\begin{equation}
    \mathbf{M} = \nabla \times (\mathbf{r} \psi), \qquad \mathbf{N} = \frac{\nabla \times \mathbf{M}}{k},
\end{equation}
兩者皆滿足向量亥姆霍茲方程、散度為零，且互為旋度，因此可用來展開電磁場。$\psi$ 在球座標下分離變數後，角向解為連帶勒讓德函數 $P_n^{\ell}(\cos\theta)$ 與 $\cos\ell\phi$、$\sin\ell\phi$，徑向解為球貝索函數 $z_n(kr)$：球內與入射場取在原點有限的 $j_n$，散射場取向外傳播的球漢克爾函數 $h_n^{(1)}$。平面波入射時只用到 $\ell=1$ 的項，定義角函數
\begin{equation}
    \pi_n = \frac{P_n^1\left(\cos\theta\right)}{\sin\theta}, \qquad \tau_n = \frac{dP_n^1\left(\cos\theta\right)}{d\theta} .
\end{equation}
這些函數在球面上彼此正交，所以展開係數可以由內積直接求出\parencite{bohren1983Absorption}。

\subsubsection*{(四) 場的級數展開}
沿 $+z$ 傳播、沿 $x$ 方向偏振的平面波，由正交關係可得只剩 $\mathbf{M}_{o1n}$ 與 $\mathbf{N}_{e1n}$ 兩組項：
\begin{equation}
    \mathbf{E}_i = \sum_{n=1}^{\infty} E_n \left( \mathbf{M}^{(1)}_{o1n} - i \mathbf{N}^{(1)}_{e1n} \right), \qquad E_n = i^n E_0 \frac{2n+1}{n(n+1)} .
\end{equation}
散射場與球內場的角向依賴和入射場相同，只差待定係數：
\begin{equation}
    \mathbf{E}_s = \sum_{n=1}^{\infty} E_n \left( i a_n \mathbf{N}^{(3)}_{e1n} - b_n \mathbf{M}^{(3)}_{o1n} \right), \qquad
    \mathbf{E}_1 = \sum_{n=1}^{\infty} E_n \left( c_n \mathbf{M}^{(1)}_{o1n} - i d_n \mathbf{N}^{(1)}_{e1n} \right) .
\end{equation}
上標 $(1)$、$(3)$ 分別代表徑向函數取 $j_n$ 與 $h_n^{(1)}$。

\subsubsection*{(五) 散射係數與截面}
把各場代入 (二) 的邊界條件，令尺寸參數 $x = ka$、相對折射率 $m = k_1/k$，並以 Riccati--Bessel 函數 $\psi_n(\rho) = \rho\, j_n(\rho)$、$\xi_n(\rho) = \rho\, h_n^{(1)}(\rho)$ 表示，解得米氏散射係數
\begin{align}
    a_n &= \frac{m\, \psi_n(mx)\, \psi_n^{\prime}(x) - \psi_n(x)\, \psi_n^{\prime}(mx)}{m\, \psi_n(mx)\, \xi_n^{\prime}(x) - \xi_n(x)\, \psi_n^{\prime}(mx)}, \\
    b_n &= \frac{\psi_n(mx)\, \psi_n^{\prime}(x) - m\, \psi_n(x)\, \psi_n^{\prime}(mx)}{\psi_n(mx)\, \xi_n^{\prime}(x) - m\, \xi_n(x)\, \psi_n^{\prime}(mx)} ,
\end{align}
$a_n$、$b_n$ 分別對應第 $n$ 階的電型與磁型多極輻射。將遠場的散射功率對全球面積分，再除以入射強度與幾何截面 $\pi a^2$，得散射效率；由光學定理可得消光效率：
\begin{equation}
    \begin{aligned}
    Q_{sca} &= \frac{2}{x^2} \sum_{n=1}^{\infty}(2n+1)\left(|a_n|^2+|b_n|^2\right), \\
    Q_{ext} &= \frac{2}{x^2} \sum_{n=1}^{\infty} (2n+1)\, \mathrm{Re}\left\{ a_n + b_n \right\} .
    \end{aligned}
    \label{eq:Qsca_proof}
\end{equation}
級數在 $n$ 略大於 $x$ 之後迅速收斂，數值計算通常取 $n_{max} \approx x + 4x^{1/3} + 2$\parencite{bohren1983Absorption}。
